\begin{itemize}
    \item
    New users are restricted to a total of 32 cores: write to \url{rt-ex-hpc@encs.concordia.ca}
    if you need more temporarily (192 is the maximum, or, 6 jobs of 32 cores each).
    
    \item
    Batch job sessions are a maximum of one week in length (only 24 hours, though,
    for interactive jobs, see \xs{sect:interactive-jobs}).
    
    \item
    Scripts can live in your NFS-provided home, but any substantial data need
    to be in your cluster-specific directory
    (located at \verb+/speed-scratch/<ENCSusername>/+).
    
    NFS is great for acute activity, but is not ideal for chronic activity.
    Any data that a job will 
    read more than once should be copied at the start to the scratch disk of a 
    compute node using \api{\$TMPDIR} (and, perhaps, \api{\$SLURM\_SUBMIT\_DIR}), 
    any intermediary job data should be produced in \api{\$TMPDIR}, and once a 
    job is near to finishing, those data should be copied to your NFS-mounted 
    home (or other NFS-mounted space) from \api{\$TMPDIR} (to, perhaps,
    \api{\$SLURM\_SUBMIT\_DIR}). In other words, IO-intensive operations should be effected 
    locally whenever possible, saving network activity for the start and end of 
    jobs. 
    
    \item
    Your current resource allocation is based upon past usage, which is an 
    amalgamation of approximately one week's worth of past wallclock (i.e., time 
    spent on the node(s)) and compute activity (on the node(s)).
    
    \item
    Jobs should NEVER be run outside of the province of the scheduler.
    Repeat offenders risk loss of cluster access. 
    
\end{itemize}